\subsection*{Solutions}

\begin{solution}
  for \textbf{Exercise} \ref{exc:which_of_the_following_represent_a_set}
  \begin{enumerate}
    \label{enum:sets_and_numbers_example_solution}

    \item This represents a set since it's well defined. We all know what it
      means to be registered for a class.
    \item This does not represent a set since it's not well defined. There are
      many different interpretations of what it means to be a good student
      (get an \textbf{A} or pass the class or attend class or avoid falling
      asleep in class or don't cause trouble in class)
  \end{enumerate}
\end{solution}

\begin{solution}
  for \textbf{Exercise} \ref{exc:simplify_the_following_expressions}

  \begin{itemize}
    \label{item:simplify_the_following_expressions_solutions}

    \item $(-4, \infty) \cup [-8, 3] = [-8, \infty)$
    \item $(-4, \infty) \cup (-\infty, 2] = (-\infty, \infty) = \mathbb{R}$
    \item $(-4, \infty) \cap (-\infty, 2] = (-4, 2]$
    \item $(-4, \infty) \cap [-10, -5] = \emptyset$
  \end{itemize}
\end{solution}

\begin{solution}
  for \textbf{Exercise} \ref{exc:convert_8_radians_into_degrees_and_8_into_radians}

  Since $2\pi = 360^{\circ} \qquad \frac{2\pi}{360^{\circ}} = 1 = \frac{360^{\circ}}{2\pi}$.
  We're trying to cancel out the radians. So, to do this, we will need to have
  the radians on the bottom, and the degrees on the top. After we do the
  multiplication, we will be left with the degrees.

  \begin{align*}
    8\rad \times \left(\frac{360^{\circ}}{2\pi \rad}\right) &=
      \frac{8 \times 360^{\circ}}{2\pi} \\
    &= \frac{1440^{\circ}}{\pi} \\
    &\approx 458.4
  .\end{align*}

  Now, let's convert $8^{\circ}$ into $8$ radians.

  \begin{align*}
    8^{\circ} \left(\frac{\pi \rad}{180^{\circ}}\right) &=
      \frac{8\pi}{180} \rad \\
    &= \frac{\pi}{45} \rad \\
  .\end{align*}
\end{solution}

\begin{solution}
  for \textbf{Exercise}
  \ref{exc:what_is_the_arc_length_spanned_by_a_40_degree_angle_on_a_circle_of_radius_30_meters}

  Convert $40^{\circ}$ into radians:

  \begin{align*}
    40^{\circ} \times \left(\frac{\pi}{180^{\circ}}\right) &=
      \frac{40\pi}{180} \rad \\
    &= \frac{4\pi}{18} \\
    &= \frac{2\pi}{9} \\
    &\approx 0.7 \\
  .\end{align*}

  Now we're ready to compute the arc length:

  \begin{align*}
    s &= r \times \mid \theta \mid \\
      &= (30m) \times (\frac{2\pi}{9}) \\
      &= \frac{20\pi}{3}m
  .\end{align*}
\end{solution}

\begin{solution}
  for \textbf{Exercise}
  \ref{exc:put_it_all_together}

  \begin{figure}[H]
    \centering
    \incfig{ferris-wheel-diagram}
    \caption{Ferris Wheel Diagram}
    \label{fig:ferris_wheel_diagram_pdf}
  \end{figure} 
\end{solution}
